\section{Extension: do published future-impact studies point the same way?}\label{sec:futurelit}
\paragraph{Purpose.} Section~\ref{sec:checks} compared our results with \litN\ published findings for a few crops. This
section asks, for \emph{every} one of the nine crops, what peer-reviewed future-impact research exists, what it
concludes, and whether our conclusion on the same crop and place points the same way. Coverage is complete when every
crop is accounted for, including crops for which no suitable study was found. It is not a score of accuracy.

\subsection{Method}
\begin{itemize}
\item \textbf{Studies.} \litNRows\ study--crop rows from \litNStudies\ publications, searched for each crop with
priority to panel countries and to the 2040s--2060s. \litNFull\ rows were read in full text or supplement and
\litNAbstract\ at abstract level only. Each row records full citation, DOI, read level, locator, period, scenario,
climate models, CO$_2$ and adaptation treatment, irrigation assumption, outcome measure, and a paraphrased conclusion
(\file{docs/evidence/future\_literature\_sources.csv}; quotations in \file{docs/evidence/future\_literature\_sources.md}). Every DOI was
resolved at Crossref.
\item \textbf{Our side.} For the panel countries a study covers, the direction of our conditional yield change
(literature damage accounting, with the present-relative variant and the empirical cross-check), the direction of the
F4 change and whether heat or water drives it are computed from the saved results
(\file{code/30\_future\_literature.py}). Where a study covers none of our panel countries, our whole panel is shown and
flagged.
\item \textbf{Comparison.} Directions, mechanisms and conclusions only. Yield percentages from different metrics,
scenarios and periods are never set side by side. A verdict (consistent, partly consistent, different, not directly
comparable) is a judgement that the direction and stated mechanism of the paper's conclusion do or do not match ours.
The verdicts were assigned in this audit and have not been reviewed by teammates.
\item \textbf{Differences are recorded, not forced.} Studies differ from our setup in region, period, scenario,
climate models, irrigation, adaptation, CO$_2$ treatment and outcome. \litNCoTwoIncl\ of the \litNRows\ rows include
CO$_2$ fertilisation (alone or with a no-CO$_2$ run) and \litNAdaptIncl\ include adaptation; our index includes
neither.
\item \textbf{Shared inputs.} Three rows come from \citet{deryng2014}, whose thresholds are our wheat, maize and soybean
parameters, so agreement with them is not independent evidence. \citet{zhao2017} is used as a benchmark elsewhere
(Section~\ref{sec:checks}).
\end{itemize}

\begin{table}[H]\centering\scriptsize
\caption[Future-literature coverage by crop]{Coverage of future-impact literature by crop. ``Full / abstract'' counts
rows by read level. Verdict columns: consistent / partly consistent / different / not directly comparable. The last
column lists panel countries for which no verified future study was found.}\label{tab:litcoverage}
\adjustbox{max width=\linewidth}{\begin{tabular}{@{}lrrrrlp{4.4cm}@{}}
\toprule
Crop & Studies & Rows & Full / abstract & Panel countries covered & Cons. / partly / diff. / n.c. & Panel countries without a verified study \\
\midrule
Wheat & 6 & 6 & 4 / 2 & 4 of 6 & 4 / 2 / 0 / 0 & Egypt, Brazil \\
Rice & 4 & 4 & 3 / 1 & 4 of 7 & 0 / 4 / 0 / 0 & USA, Nigeria, Italy \\
Maize & 6 & 6 & 6 / 0 & 5 of 6 & 1 / 5 / 0 / 0 & South Africa \\
Chickpea & 2 & 3 & 1 / 2 & 2 of 7 & 2 / 1 / 0 / 0 & Australia, T\"urkiye, Russia, USA, China \\
Lentil & 3 & 3 & 1 / 2 & 2 of 7 & 0 / 2 / 1 / 0 & Australia, USA, Russia, China, Ethiopia \\
Common bean & 4 & 4 & 2 / 2 & 2 of 7 & 2 / 1 / 1 / 0 & India, Myanmar, China, USA, Belarus \\
Soybean & 5 & 5 & 5 / 0 & 6 of 7 & 0 / 4 / 1 / 0 & South Africa \\
Groundnut & 4 & 4 & 3 / 1 & 2 of 5 & 1 / 1 / 2 / 0 & China, USA, Brazil \\
Rapeseed-mustard & 4 & 4 & 1 / 3 & 3 of 7 & 2 / 1 / 0 / 1 & China, USA, South Africa, Brazil \\
\bottomrule
\end{tabular}
}
\end{table}

Of the \litNPairs\ crop--country pairs of the panel, \litNPairsNoStudy\ have no verified future study at all. Counting
the verdicts (\litNCons\ consistent, \litNPartly\ partly consistent, \litNDiff\ different, \litNNC\ not directly
comparable) says only how the evidence is distributed; it does not measure the model's accuracy.

\subsection{Crop by crop}
\paragraph{Wheat.} Wheat has the strongest and most consistent support. \citet{zhao2017} find a negative temperature
effect with every method, larger for India and Russia than for the USA and smallest for China. Our conditional yield
change is negative in all four countries and also smallest in China. For India, \citet{nareshkumar2014wheat} project a
6--23\% loss by 2050 even with CO$_2$ fertilisation, and \citet{chandran2023} find a 7--9\% loss from a shortened crop
cycle; both agree with a heat-driven loss for India. \citet{asseng2015} give about $-6\%$ per $^\circ$C
(abstract only). Two studies only partly agree. \citet{jagermeyr2021} project global wheat gains with CO$_2$
fertilisation, losses in South Asian spring wheat and gains in the North China Plain and northern winter wheat: our
losses for China and the USA conflict with that pattern. \citet{deryng2014} find the same loss as ours without CO$_2$
but gains with it. In Russia the literature stresses temperature while our change is mainly water-driven, a mechanism
difference noted earlier \citep{lobell2011}.

\paragraph{Rice.} Rice is the weakest spot of our setup. \citet{zhao2017} find large losses for India (about
$-6.6\%$ per $^\circ$C) but small averaged impacts for China and Bangladesh. Ours agrees for China and Bangladesh and
gives about zero for India on the main route, because monthly flowering temperature rarely exceeds the 35\,$^\circ$C
threshold; our empirical cross-check ($-11.9\%$) agrees with them. \citet{soora2013} project a 7\% loss for irrigated rice
by 2050 and only a marginal loss for rainfed rice; irrigated area carries no water stress in our index, so we cannot
reproduce the irrigated result. \citet{jagermeyr2021} and \citet{chandran2023} find gains with CO$_2$ fertilisation and
a result spanning zero without it. A low rice score here therefore reflects what the model cannot see (hot days,
irrigation) as much as low exposure.

\paragraph{Maize.} Studies agree on losses \citep{zhao2017,jagermeyr2021,tigchelaar2018,zilli2020}. Our signs match
for the USA, Brazil and Ukraine on the main route. For China and India the main route gives small gains, because our
maize score is water-dominated and the future is wetter there, whereas the temperature-only studies show losses; our
empirical cross-check matches them. \citet{tigchelaar2018} use the method closest to ours (delta change, no CO$_2$,
no adaptation) and their country losses are mirrored by our empirical route in all five countries.
\citet{jagermeyr2021} note that MRI-ESM2, our only climate model, gives the least pessimistic crop impacts of the
models they used, so our losses are likely at the mild end of the range.

\paragraph{Soybean.} The global signal is weak (confidence intervals include zero in \citealt{zhao2017}; an interquartile
range from $-21$ to $+17\%$ in \citealt{jagermeyr2021}), with losses in the USA and gains in parts of China, which ours
reproduces. For Indian soybean \citet{walikar2018} report small declines from a shorter crop duration, while ours is
about zero to slightly positive, a difference we cannot attribute to a single cause. \citet{zilli2020} find Brazilian
losses, concentrated in the northern Cerrado, where ours is near zero on the main route.

\paragraph{Chickpea.} Only \citet{singh2014chickpea} and \citet{abera2026} offer future simulations, none for five of
the seven panel countries. They find gains for Ethiopia, as ours does, and site-dependent changes in India that
include CO$_2$ effects. Neither shows a uniform large loss from climate alone; our main route is near zero for India
and our empirical route negative.

\paragraph{Lentil.} \citet{qian2026} project Prairie lentil gains of about 20--28\% in the 2050s from CO$_2$ and
warmth with adjusted planting dates (abstract only); ours is a small loss, a difference that coincides with the
CO$_2$ and adaptation treatment but cannot be separated from other differences. For India, the West Bengal studies
\citep{chandran2022,chandran2023} give losses of 12--31\% attributed to a shortened reproductive phase, where ours is
small and water-driven. These come from one site and one research group.

\paragraph{Common bean.} For Tanzania, suitability studies \citep{rippke2016,hummel2018} find beans the most exposed
staple, from heat in the north-west and rainfall decline in the east, and our large Tanzanian loss with equal heat and
water shares matches both direction and mechanism. For Brazil, \citet{antolin2021} project gains (with CO$_2$) where ours
is a large heat-driven loss, and \citet{heinemann2017} stress drought frequency, a mechanism ours shows less clearly.
No study covers the Asian beans that dominate the FAOSTAT item in India and Myanmar (Section~\ref{sec:parevidence}).

\paragraph{Groundnut.} \citet{kadiyala2021} find gains where rainfall rises and losses where warming hastens maturity,
a mechanism that matches our India result, a small gain from a wetter season. The West Bengal studies
\citep{chandran2022,chandran2023} find losses of up to 48\% at one site; ours differs. For Nigeria,
\citet{rippke2016} find limited transformation pressure, which ours does not contradict.

\paragraph{Rapeseed-mustard.} Studies of Indian mustard \citep{nareshkumar2016mustard,boomiraj2010} find losses driven
by maximum temperature, which matches our India result in direction and mechanism. This agreement cannot be read as
support for our rapeseed parameters, because the heat slope is unsupported (Section~\ref{sec:parevidence}). For Canada,
\citet{qian2019} find losses with CO$_2$ fixed and gains with CO$_2$; ours matches the first. The European study
\citep{pullens2019} reports frost and drought risk types rather than yield and is not directly comparable.

\subsection{What the comparison shows and what it cannot}
\begin{itemize}
\item \textbf{Supportive.} Direction and mechanism agree for wheat in India, bean in Tanzania, mustard in India and
chickpea in Ethiopia, and the weak soybean signal is also seen elsewhere.
\item \textbf{Challenging.} \litNChallenge\ rows report gains or mixed outcomes, mostly where CO$_2$ fertilisation or
adaptation is included: Canadian lentil, Brazilian bean, wheat in China and the USA, rice in South Asia, and
groundnut under one of two climate models. Where studies find losses and ours does not (Indian soybean and groundnut,
Indian rice), the cause is more often missing mechanisms: season shortening, reproductive-phase damage, hot days and
irrigation.
\item \textbf{A low score can hide mechanisms.} Rice, soybean, lentil and groundnut scores are low partly because the
index sees only monthly heat above a threshold and rainfed deficit. Low climate-explained yield variability
\citep{ray2015} is a different outcome from low vulnerability and is not treated as evidence for it.
\item \textbf{National averages hide winners and losers.} Several studies find opposite signs within one country
(the North China Plain against South Asia, or Indian basins against each other); our values are national.
\item \textbf{Limits.} Differences in CO$_2$, adaptation, models, scenarios and outcome measures coincide, so a
disagreement cannot be assigned to one cause. \litNAbstract\ of the \litNRows\ rows are read at abstract level only,
and the West Bengal rows come from one site. Agreement with a study that supplies our inputs is not independent validation. Leads that could not be used
(Moura da Silva 2021 for Brazilian soybean, Boomiraj 2013, Singh 2014 groundnut) are listed in
\file{docs/evidence/future\_literature\_sources.md}.
\end{itemize}
Full row-level comparison: Table~\ref{tab:litrows} in Appendix~\ref{app:ext}.
